\chapter*{概要}
\addcontentsline{toc}{chapter}{概要}

\SP{} では、LLM の決算発表文の読み取りは発表の当日に株価に入りきり、翌日以降のドリフトは予測しなかった（実験 2・3・8）。そこで、注目の低い小型株（浮動株時価総額で 1,001〜3,000 位、1,000 社）で、読み取りが発表の後 2〜60 取引日のドリフトを予測するかを、事前登録して 1 回だけ確かめた。

\begin{keybox}[要点]
\begin{itemize}
\item \textbf{主な検定は支持された}。ChronoGPT（2015 年版）のベクトルに、2 日の反応を正解にしたリッジを当てた点数が 1 標準偏差高い発表は、2 日の反応・SUE などを統制した後、2〜60 日の IWM との超過リターンが \V{Pcoef} bp 高かった（$t$ = \V{Pt}、置換の片側 $p$ \V{Pperm}、発表 \V{Pn} 件）。2021 年以降だけでも \V{Rcoef} bp（$t$ = \V{Rt}）。
\item \textbf{gpt-4o-mini}（2024 年以降）は、学習なしで \V{gptcoefS} bp（$t$ = \V{gpttS}）、重みを学習して \V{gptridgecoefS} bp（$t$ = \V{gptridgetS}）と、さらに強かった。学習なしの辞書は負、FinBERT はゼロだった。
\item \textbf{ドリフトは主に 22〜60 日目}に出ていた（ChronoGPT）。過去 1 年のモメンタムを統制に加えても消えず、規模の大小のどちらの半分でも見られた（事前登録の後の探索）。
\item \textbf{取引の成績にはならなかった}。浮動株時価総額で加重した上位 20\% の月次のポートフォリオは、片道 0.3\% の費用の後で、範囲の全社を月 \V{Mexb} 上回ったが、有意ではなかった（$t$ = \V{Mtb}、\V{Mmonths} か月）。
\item \textbf{BLM の資産の推移}（追記 1、検定はしない）：ChronoGPT＋リッジの見通しの BLM は、2015 年末からの四半期ごとのリバランスで、最終の値が \V{B2015quarterchronoridge}、比べる相手（範囲の全社、浮動株時価総額で加重）が \V{B2015quarterband}（年率の差 \V{B2015quarterchronoridgeex}）だった。半年・1 年ごとでも、2023 年末からでも上回った。ただし差の多くは 2023 年以降に開き、資産の水準は生存者バイアスで高く出ている。

\item \textbf{注意}：標本は今も上場している会社だけである（生存者バイアス）。収益性・アナリストの予想の修正は統制していない。結果は 1 回の実行のものである。
\end{itemize}
\end{keybox}
